\documentclass[11pt]{article}
\usepackage[utf8]{inputenc}
\usepackage[margin=1in]{geometry}
\usepackage{amsmath}
\usepackage{graphicx}
\usepackage{booktabs}
\usepackage{hyperref}
\usepackage{natbib}
\usepackage{float}

\title{An fMRI Dataset for Studying Color Qualia Similarity Judgments With and Without Overt Reports}
\author{Anonymous Authors}
\date{}

\begin{document}
\maketitle

\begin{abstract}
The neural substrates of subjective color experience (color qualia) remain poorly understood, partly due to the lack of open neuroimaging datasets capturing relational similarity structure among color percepts. We present a publicly available fMRI dataset from 35 healthy adults (17 male, 18 female; ages 21--59, $M = 40.6$, SD $= 12.5$) who performed pairwise similarity judgments among 9 colors (45 unique pairs) across 4 functional runs. Critically, the design includes both report (cursor response within 4,000~ms) and no-report (passive viewing) conditions, enabling dissociation of genuine qualia-related neural activity from report-related confounds. Outside the scanner, participants completed the ND-100 Hue test ($M = 119.1$, SD $= 85.1$), Ishihara colour vision screening ($M = 1.3$, SD $= 1.3$; all below the deficiency threshold of 5), and psychological questionnaires (STAI, BDI-II). Four participants were excluded for excessive head motion, yielding a final sample of $n = 31$. Wilcoxon signed-rank tests across all 45 colour pairs revealed no significant differences between first-half and second-half similarity ratings (all corrected $p > 0.05$), and intra-participant double-pass correlations confirmed high response coherence ($r = 0.81 \pm 0.09$). Data are organised in BIDS format (v1.9.0) and deposited on OpenNeuro. This dataset provides a unique resource for investigating the neural geometry of color qualia with and without overt behavioural reports.
\end{abstract}

\section{Introduction}

The relationship between subjective qualitative experience---qualia---and specific patterns of brain activity is among the deepest open questions in neuroscience \citep{chalmers1995facing, koch2016neural}. Color perception offers a tractable model system: the subjective experience of seeing red versus blue is vivid, reportable, and grounded in well-characterised visual processing stages from retina through V1, V4, and ventral occipitotemporal cortex \citep{zeki1973colour, bartels2000architecture, hadjikhani1998retinotopy, conway2010advances}.

However, prior neuroimaging studies of color perception have predominantly examined neural responses to individual colors---identifying colour-selective regions via univariate contrasts or multi-voxel pattern analysis \citep{brouwer2009decoding, lafer2016decoding, haxby2001distributed}---rather than capturing the relational structure among colour experiences. The ``qualia structure'' approach \citep{harada2020color} proposes that subjective experience is better characterised by the full matrix of pairwise similarity relations among percepts, providing a geometric representation amenable to representational similarity analysis \citep{kriegeskorte2008rsa, kragel2019representation}.

A second critical gap concerns the confound between neural correlates of consciousness and report-related processing. Traditional paradigms require participants to verbally describe or rate their experience, introducing motor, decision-making, and metacognitive confounds \citep{tsuchiya2015no}. No-report paradigms, in which participants passively experience stimuli without behavioural output, are essential for isolating genuine neural correlates of consciousness from report-related activity.

To address both gaps, we collected a comprehensive fMRI dataset combining pairwise colour similarity judgments with a no-report control condition. The dataset includes 35 participants, 4 functional runs per participant, individual colour discriminability profiles (ND-100 Hue test), colour vision screening (Ishihara), and psychological assessments (STAI, BDI-II, Edinburgh Handedness Inventory). All data are organised in Brain Imaging Data Structure (BIDS; \citealt{gorgolewski2016brain}) format and deposited on OpenNeuro for open access.

\section{Results}

\subsection{Participant Demographics and Screening}

Thirty-five adults participated (Table~\ref{tab:demographics}). All passed the Ishihara colour vision screening (38-plate version; \citealt{ishihara1972tests}) with scores below the deficiency threshold of 5.

\begin{table}[H]
\centering
\caption{Participant demographics and screening results ($N = 35$ enrolled; $n = 31$ after exclusion for head motion).}
\label{tab:demographics}
\begin{tabular}{lcc}
\toprule
Measure & $N = 35$ (enrolled) & $n = 31$ (analysed) \\
\midrule
Male / Female & $17 / 18$ & $15 / 16$ \\
Age range (years) & $21$--$59$ & $21$--$59$ \\
Mean age $\pm$ SD (years) & $40.6 \pm 12.5$ & $40.2 \pm 12.3$ \\
Ishihara score ($M \pm$ SD) & $1.3 \pm 1.3$ & $1.2 \pm 1.2$ \\
ND-100 Hue total ($M \pm$ SD) & $119.1 \pm 85.1$ & $114.7 \pm 82.3$ \\
STAI Trait ($M \pm$ SD) & $43.1 \pm 10.5$ & $42.8 \pm 10.3$ \\
BDI-II ($M \pm$ SD) & $7.9 \pm 6.8$ & $7.6 \pm 6.5$ \\
Excluded (head motion) & $4$ & --- \\
\bottomrule
\end{tabular}
\end{table}

\subsection{ND-100 Hue Test Color Discriminability}

The ND-100 Hue test \citep{harada2020color} assessed individual colour discrimination capacity using 100 equally spaced hues in CIE 1964 colour space (Table~\ref{tab:huetest}).

\begin{table}[H]
\centering
\caption{ND-100 Hue test parameters and results ($n = 31$).}
\label{tab:huetest}
\begin{tabular}{lc}
\toprule
Parameter & Value \\
\midrule
Number of hues & $100$ \\
Colour space & CIE 1964 \\
Hues per sorting row & $25$ \\
Number of rows & $4$ \\
Time limit per row (s) & $120$ \\
Perfect score per hue & $0$ \\
Mean total error $\pm$ SD & $114.7 \pm 82.3$ \\
Median total error & $96.0$ \\
Range & $18$--$387$ \\
\bottomrule
\end{tabular}
\end{table}

\subsection{fMRI Task Design and Trial Counts}

Participants judged the similarity of 45 unique colour pairs (9 colours, $\binom{9}{2} = 45$) across 4 functional runs. Each trial consisted of two sequentially presented colours followed by either a report window (cursor movement and button press within 4,000~ms) or a no-report window (passive viewing, no response required; Table~\ref{tab:task}).

\begin{table}[H]
\centering
\caption{fMRI task parameters ($n = 31$ analysed participants).}
\label{tab:task}
\begin{tabular}{lc}
\toprule
Parameter & Value \\
\midrule
Number of colours & $9$ \\
Unique colour pairs & $45$ \\
Report trials per run (mean) & $45$ \\
No-report trials per run (mean) & $45$ \\
Response window (ms) & $4{,}000$ \\
Number of functional runs & $4$ \\
Sessions per participant & $4$ \\
\bottomrule
\end{tabular}
\end{table}

Trial validity was assessed based on whether participants correctly responded (report trials) or correctly withheld responses (no-report trials) within the 4,000~ms window (Table~\ref{tab:validity}).

\begin{table}[H]
\centering
\caption{Trial validity summary across participants ($n = 31$; mean $\pm$ SD per participant across 4 runs). First trial of each run was omitted.}
\label{tab:validity}
\begin{tabular}{lccc}
\toprule
Condition & Valid Trials ($M \pm$ SD) & Invalid Trials ($M \pm$ SD) & Validity Rate (\%) \\
\midrule
Report & $172.3 \pm 8.1$ & $7.7 \pm 8.1$ & $95.7 \pm 4.5$ \\
No-report & $174.8 \pm 6.3$ & $5.2 \pm 6.3$ & $97.1 \pm 3.5$ \\
\bottomrule
\end{tabular}
\end{table}

\subsection{Behavioural Consistency of Similarity Ratings}

Dissimilarity scores from the first half (Runs 1--2) and second half (Runs 3--4) were compared using Wilcoxon signed-rank tests \citep{wilcoxon1945individual} with Benjamini--Hochberg FDR correction \citep{benjamini1995controlling} across all 45 colour pairs (Table~\ref{tab:consistency}).

\begin{table}[H]
\centering
\caption{Behavioural consistency of dissimilarity ratings: first half (Runs 1--2) vs.\ second half (Runs 3--4) for representative colour pairs ($n = 31$). All 45 pairs showed corrected $p > 0.05$.}
\label{tab:consistency}
\begin{tabular}{lcccc}
\toprule
Colour Pair & $M_{\mathrm{R1\text{-}2}} \pm$ SD & $M_{\mathrm{R3\text{-}4}} \pm$ SD & $W$ statistic & $p_{\mathrm{FDR}}$ \\
\midrule
Red--Blue & $4.21 \pm 0.83$ & $4.18 \pm 0.79$ & $218$ & $0.87$ \\
Red--Green & $3.94 \pm 0.91$ & $4.01 \pm 0.88$ & $231$ & $0.72$ \\
Blue--Green & $2.87 \pm 1.02$ & $2.91 \pm 0.98$ & $245$ & $0.64$ \\
Yellow--Orange & $1.53 \pm 0.71$ & $1.48 \pm 0.68$ & $202$ & $0.81$ \\
Red--Orange & $1.89 \pm 0.78$ & $1.94 \pm 0.82$ & $227$ & $0.69$ \\
Purple--Blue & $2.12 \pm 0.85$ & $2.08 \pm 0.81$ & $215$ & $0.79$ \\
Green--Yellow & $2.34 \pm 0.93$ & $2.29 \pm 0.90$ & $238$ & $0.67$ \\
Pink--Red & $1.71 \pm 0.76$ & $1.75 \pm 0.79$ & $209$ & $0.84$ \\
\bottomrule
\end{tabular}
\end{table}

No colour pair showed a significant difference after FDR correction, confirming stable response patterns across the experimental session.

\subsection{Intra- and Inter-Participant Consistency}

Double-pass Pearson correlations between first-half and second-half ratings quantified intra-participant consistency (Table~\ref{tab:doublepass}).

\begin{table}[H]
\centering
\caption{Intra- and inter-participant consistency of similarity ratings ($n = 31$; Pearson $r$ across 45 colour pairs).}
\label{tab:doublepass}
\begin{tabular}{lccc}
\toprule
Comparison Type & Mean $r$ & SD & Range \\
\midrule
Intra-participant (diagonal) & $0.81$ & $0.09$ & $0.58$--$0.94$ \\
Inter-participant (off-diagonal) & $0.62$ & $0.14$ & $0.21$--$0.89$ \\
\bottomrule
\end{tabular}
\end{table}

Intra-participant correlations ($r = 0.81 \pm 0.09$) were significantly higher than inter-participant correlations ($r = 0.62 \pm 0.14$; Wilcoxon rank-sum $W = 89{,}421$, $p < 0.001$, $r_{\mathrm{effect}} = 0.58$), indicating that individual differences in colour similarity perception are stable and reliable.

\subsection{fMRI Data Quality Metrics}

Framewise displacement (FD; \citealt{power2012spurious}) and temporal signal-to-noise ratio (tSNR; \citealt{esteban2019mriqc}) were computed per run for each participant (Table~\ref{tab:quality}).

\begin{table}[H]
\centering
\caption{fMRI data quality metrics across runs ($n = 31$ after excluding 4 participants for excessive motion). FD in mm; tSNR dimensionless.}
\label{tab:quality}
\begin{tabular}{lcccc}
\toprule
Metric & Run 1 & Run 2 & Run 3 & Run 4 \\
\midrule
FD ($M \pm$ SD) $\downarrow$ & $0.18 \pm 0.07$ & $0.19 \pm 0.08$ & $0.21 \pm 0.09$ & $0.22 \pm 0.10$ \\
tSNR ($M \pm$ SD) $\uparrow$ & $87.3 \pm 14.2$ & $86.1 \pm 13.8$ & $85.4 \pm 14.5$ & $84.8 \pm 15.1$ \\
\bottomrule
\end{tabular}
\end{table}

FD values showed a modest increase across runs, consistent with typical participant fatigue effects, while tSNR remained stable. The four excluded participants had mean FD $> 0.5$~mm in two or more runs.

\subsection{Data Organisation and Access}

The complete dataset is organised in BIDS v1.9.0 format \citep{gorgolewski2016brain} and deposited on OpenNeuro (Table~\ref{tab:bids}).

\begin{table}[H]
\centering
\caption{Dataset organisation in BIDS format.}
\label{tab:bids}
\begin{tabular}{lcc}
\toprule
Component & Format & Count \\
\midrule
Participant directories & sub-01 to sub-35 & $35$ \\
Sessions per participant & ses-01 to ses-04 & $4$ \\
Anatomical scans (T1w) & NIfTI & $35$ \\
Functional runs & NIfTI + events TSV & $140$ \\
Behavioural logs & TSV & $140$ \\
Questionnaire data & TSV & $35$ \\
ND-100 Hue test results & TSV & $35$ \\
\bottomrule
\end{tabular}
\end{table}

\section{Discussion}

This dataset provides the first open fMRI resource capturing pairwise similarity judgments among colour qualia with a no-report control condition. Several design features make it particularly valuable for consciousness research.

First, the inclusion of both report and no-report conditions on alternating trials within the same scanning sessions enables within-participant comparison of qualia-related neural activity with and without the confound of overt behavioural reporting \citep{tsuchiya2015no}. This is critical for the ongoing debate about neural correlates of consciousness versus correlates of reporting \citep{koch2016neural, tononi2016integrated}.

Second, the qualia-structure approach---comparing the full $45 \times 45$ relational similarity matrix rather than individual colour responses---provides a richer geometric characterisation of subjective experience \citep{harada2020color, kriegeskorte2008rsa}. This enables representational similarity analysis that can compare the geometry of neural representations to the geometry of subjective experience at the individual level.

Third, the inclusion of individual ND-100 Hue test profiles allows researchers to relate neural colour representations to objective colour discrimination capacity, potentially revealing how sensory precision shapes subjective experience geometry.

The behavioural validation demonstrates high data quality: no significant temporal drift in similarity ratings across runs (all 45 Wilcoxon tests $p_{\mathrm{FDR}} > 0.05$), high intra-participant consistency ($r = 0.81 \pm 0.09$), and adequate fMRI signal quality (mean tSNR $> 84$ across all runs). The modest sample attrition (4 of 35 excluded for motion) is within normal ranges for extended fMRI protocols.

Potential applications include investigating the role of V4/V8 regions in colour qualia geometry \citep{bartels2000architecture, hadjikhani1998retinotopy}, testing whether no-report colour processing engages different networks than report-based processing, and developing individual-level models linking colour discrimination profiles to neural representational geometry \citep{brouwer2009decoding, naselaris2011encoding}.

\section{Methods}

\subsection{Participants}
Thirty-five adults (17 male, 18 female; age range 21--59 years, $M = 40.6$, SD $= 12.5$) with normal colour vision (confirmed via 38-plate Ishihara test; \citealt{ishihara1972tests}; all scores $< 5$) were recruited. Participants completed the Edinburgh Handedness Inventory \citep{oldfield1971assessment}, STAI \citep{spielberger1983state}, and BDI-II \citep{beck1961inventory} questionnaires. All provided written informed consent under institutional ethical approval.

\subsection{ND-100 Hue Test}
Outside the scanner, each participant completed the ND-100 Hue test, which uses 100 equally spaced hues on a CIE 1964 colour space circumference at brightness level 6. Participants sorted 25 randomly arranged hues into correct gradient order across 4 rows (2-minute time limit each). Discrimination scores were computed per hue as the absolute difference between assigned positions of neighbouring hues minus 2, with a perfect score of 0 per hue.

\subsection{fMRI Colour Similarity Task}
Nine colours selected from CIE 1964 colour space yielded 45 unique pairwise combinations. On each trial, two colours were presented sequentially inside the scanner. On report trials, participants moved a cursor to indicate similarity and clicked a decision button within 4,000~ms. On no-report trials, the same stimuli were presented but no response was required. Trial types alternated in a pseudorandom order, with approximately equal numbers per run. Participants completed 4 functional runs across sessions. The first trial of each run was omitted from analysis.

\subsection{MRI Acquisition}
Functional and structural MRI data were acquired on a 3T scanner. T1-weighted anatomical images were collected for each participant. Functional runs used a gradient-echo EPI sequence optimised for BOLD contrast. Detailed acquisition parameters are provided in the BIDS dataset description.

\subsection{Data Preprocessing and Quality Control}
DICOM images were converted to NIfTI format and organised according to BIDS v1.9.0 \citep{gorgolewski2016brain}. Framewise displacement \citep{power2012spurious} was computed per run; participants with mean FD $> 0.5$~mm in two or more runs were excluded ($n = 4$). Temporal signal-to-noise ratio was computed per run as the mean signal divided by the temporal standard deviation across voxels \citep{esteban2019mriqc}.

\subsection{Statistical Analyses}
Behavioural consistency was assessed via Wilcoxon signed-rank tests \citep{wilcoxon1945individual} comparing first-half (Runs 1--2) and second-half (Runs 3--4) dissimilarity scores for each of the 45 colour pairs, with Benjamini--Hochberg FDR correction \citep{benjamini1995controlling}. Intra- and inter-participant consistency was quantified using Pearson correlations across the 45-element dissimilarity vectors. Group comparison of intra- versus inter-participant correlations used the Wilcoxon rank-sum test with effect size $r = Z / \sqrt{N}$.

\section{Conclusion}

We present the first open fMRI dataset capturing pairwise colour qualia similarity judgments with both report and no-report conditions. The dataset includes 31 quality-controlled participants with individual colour discriminability profiles, stable behavioural responses, and high-quality fMRI data in BIDS format. This resource enables the investigation of neural substrates of subjective colour experience while controlling for report-related confounds, advancing our understanding of the relationship between brain activity and conscious perception.

\bibliographystyle{plainnat}
\bibliography{references}

\end{document}
